\subsection{Formal Syntax of \GRFRP}
\label{sec:semantics:grfrp:syntax}

The syntax of \GRFRP is presented below and described in the following paragraphs.

\begin{center}
    \begin{bnf}(comment = {:is:})
        $F, f, x$     :is: Identifiers      :in: \idSet ;;
        $\Delta$      :is: Constants        :in: \textit{Const} ;;
        $T$           :is: Time             :in: \textit{Time} ;;
        $\mathcal{S}$ :is: Services         ::=  \service{F}{stmt^+} ;;
        $stmt$        :is: Flow Statements  ::=  \decl{x^+}{\sigma} // \import{x} // \export{x} ;;
        $\sigma$      :is: Flows            ::=  $x$ // $f(\sigma^+)$ // \onchange{\sigma} // \persist{\sigma}
        | \scan{\sigma}{T}// \sample{\sigma}{T} // \timeout{\sigma}{T}
        | \mergeop{\sigma}{\sigma} // \throttle{\sigma}{\Delta} // \timeop ;;
    \end{bnf}
\end{center}

In \GRFRP, flows are categorized as \emph{signals} (continuous values) and \emph{events} (discrete
occurrences). A signal always has a value that may evolve over time. In contrast, an event only
carries a value at the moment it occurs.

A service (\service{F}{stmt^+}) defines a flow-processing unit. Its declaration assigns it a name
$F$ and specifies a list of flow statements ($stmt^+$) describing its behavior.
%
Flow statements are of two kinds: declarations statements (\decl{x^+}{\sigma}) define the behavior
of local flows $x^+$ as flow operation $\sigma$, and interface statements (\import{x} or \export{x})
where imported flows are produced by other entities (sensors, or other services), and exported flows
are made available for other entities in the system (actuators, or other services).
%
As an example, \Cref{fig:semantics:example:aeb:code} defines the \idf{aeb} service
(\Crefrange{ex:line:aebStart}{ex:line:aebEnd}). It imports a signal \idf{speed} and events
\idf{pedestrian\_l} and \idf{pedestrian\_r}
(\Crefrange{ex:line:importSpeed}{ex:line:importPedestrians}) from sensors or other services; it
defines local flows \idf{pedestrian}, \idf{t\_pedestrian}, and \idf{strat}
(\Crefrange{ex:line:pedestrian}{ex:line:strat}); and it exports the signal \idf{strat}
(\Cref{ex:line:export}).

Flows ($\sigma$) are constructed using identifiers ($x$), applications of components
user-defined in \GRSync ($f(\sigma^+)$), and a collection of built-in operators from the \GReact
library.

The \GRSync language, used to define components, is presented in \Cref{sec:semantics:grsync}.
%
Components called in services should be \emph{reactive}, it is a crucial static property (formally defined
in \Cref{sec:semantics:grfrp:react_typing,sec:semantics:grsync:react_typing}) that
ensures a sound usage of components in asynchronous environments.
\Cref{thm:semantics:grsync:theorems:oversampling} states that a reactive component can be executed
based on changes propagation, in an event-driven manner. This execution is proved to be resilient to
\emph{oversampling}: the absence of changes from the imported flows of the service leads to an
absence of changes from its exported flows.

The \GReact library, presented in the table below, extends the \GRFRP model with flow operators
inspired by ReactiveX~\cite{reactivex}.
%
\begin{center}
    \begin{tabularx}{\textwidth}{|l|X|}\hline
        \onchange{s}         & Produces an event that emits values whenever the signal $s$ updates.                               \\\hline
        \persist{e}          & Produces a signal that updates to hold the most recent value emitted by the event $e$.             \\\hline
        \scan{s}{T}          & Produces a signal that updates with the latest value of signal $s$ every $T$ milliseconds.         \\\hline
        \sample{e}{T}        & Produces an event that emits the latest value observed from $e$ every $T$ milliseconds.            \\\hline
        \timeout{e}{T}       & Produces an event that emits a unit value if no value is emitted by $e$ within $T$ milliseconds.
        The timeout resets each time $e$ emits a value or $T$ milliseconds have passed.                                           \\\hline
        \mergeop{e_1}{e_2}   & Produces an event that emits values from either $e_1$ or $e_2$, prioritizing $e_1$ if both emit
        simultaneously.                                                                                                           \\\hline
        \throttle{s}{\Delta} & Produces a signal that updates only when the value of $s$ changes by more than threshold $\Delta$. \\\hline
        \timeop              & Produces a signal that reflects the current time.                                                  \\\hline
    \end{tabularx}
\end{center}
%
These operators enable declarative flow definitions within \GRFRP services.
%
As seen in \Cref{fig:semantics:example:aeb}, the \kwf{merge} operator is used
(\Cref{ex:line:pedestrian}) to create an event from the union of \idf{pedestrian\_l} and
\idf{pedestrian\_r}. Similarly, \kwf{timeout} (\Cref{ex:line:timeout}) triggers an event when
\idf{pedestrian} has not occurred for 2 seconds.
